\section{The Riemann integral and integral inequalities}\label{p:integration}
\blurb{Darboux sums, integrals depending on a parameter, recursions for integrals, Cauchy-Schwarz and Jensen.}

\begin{problem}
Let $f$ be a bounded function on $[0,1]$. For any partition $P = \set{x_0 < x_1 < \cdots < x_n}$ of $[0,1]$, let $\mu(P)$ denote the mesh of $P$ and
$$ R(P) = \sum_{i=1}^n (M_i - m_i)(x_i - x_{i-1}), $$
where $m_i = \inf\{f(x) : x_{i-1} \leq x \leq x_i\}$ and $M_i = \sup\{f(x) : x_{i-1} \leq x \leq x_i\}$. Prove that the following two assertions are equivalent (i.e., (i) $\implies$ (ii) and (ii) $\implies$ (i)).
\begin{enumerate}
    \item[(i)] Given $\epsilon > 0$ there exists $\delta > 0$ such that $R(P) < \epsilon$ whenever $\mu(P) < \delta$.
    \item[(ii)] There exists a real number $A$ for which given $\epsilon > 0$ there exists $\delta > 0$ such that
    $$ \abs{\sum_{i=1}^n f(t_i) (x_i - x_{i-1}) - A} < \epsilon $$
    for every partition $P$ and $x_{i-1} \leq t_i \leq x_i$, $i = 1, \ldots, n$ with $\mu(P) < \delta$.
\end{enumerate}
\begin{solution}
For any partition $P = \set{x_0 < x_1 < \cdots < x_n}$ of $[0,1]$, let
\begin{align*}
    L_P &= \sum_{i=1}^n m_i (x_i - x_{i-1}), \\
    U_P &= \sum_{i=1}^n M_i (x_i - x_{i-1}).
\end{align*}
(These are called the lower and upper Darboux sums.) Notice that
$$ R(P) = U_P - L_P $$
by definition of $R(P)$.
Since $f$ is bounded, there is an $M \geq 0$ such that $|f(x)| \leq M$ for all $x \in [0,1]$; we have that $L_P \leq M$ and $U_P \geq -M$ for any partition $P$, so we can define
\begin{align*}
    L &= \sup\{L_P \mid P \text{ is a partition of } [0,1]\}, \\
    U &= \inf\{U_P \mid P \text{ is a partition of } [0,1]\}.
\end{align*}
\begin{itemize}
    \item[] [(i) $\implies$ (ii)] For any $\epsilon > 0$, (i) implies that
    $$ U - L \leq U_P - L_P = R(P) < \epsilon $$
    for any fine enough partition $P$. Since $\epsilon$ was arbitrary, we must have $L = U$. Set $A = L = U$.
    
    Let $\epsilon > 0$. For any partition $P$, by definition of $m_i$ and $M_i$, we must have
    $$ L_P = \sum_{i=1}^n m_i (x_i - x_{i-1}) \leq \sum_{i=1}^n f(t_i) (x_i - x_{i-1}) \leq \sum_{i=1}^n M_i (x_i - x_{i-1}) = U_P. $$
    Let $\delta$ be such that $R(P) < \epsilon$ whenever $\mu(P) < \delta$, and let $P$ be such a partition. Then we find that
    $$ A - \epsilon < L_P \leq \sum_{i=1}^n f(t_i) (x_i - x_{i-1}) \leq U_P < A + \epsilon, $$
    which implies that
    $$ \abs{\sum_{i=1}^n f(t_i) (x_i - x_{i-1}) - A} < \epsilon. $$
    
    \item[] [(ii) $\implies$ (i)] Let $\epsilon > 0$ and let $\delta$ be such that
    $$ \abs{\sum_{i=1}^n f(t_i) (x_i - x_{i-1}) - A} < \frac{\epsilon}{4} $$
    for every partition $P$ and $x_{i-1} \leq t_i \leq x_i$, $i = 1, \ldots, n$ with $\mu(P) < \delta$, and let $P$ be such a partition. The infimum $m_i$ and the supremum $M_i$ need not be attained ($f$ is not assumed to be continuous), but by their definition there exist $\ell_i, u_i \in [x_{i-1}, x_i]$ such that $f(\ell_i) < m_i + \frac{\epsilon}{4}$ and $f(u_i) > M_i - \frac{\epsilon}{4}$. Since $\sum_{i=1}^n (x_i - x_{i-1}) = 1$, our assumption implies that
    \begin{align*}
        L_P &> \sum_{i=1}^n f(\ell_i) (x_i - x_{i-1}) - \frac{\epsilon}{4} > A - \frac{\epsilon}{2}, \\
        U_P &< \sum_{i=1}^n f(u_i) (x_i - x_{i-1}) + \frac{\epsilon}{4} < A + \frac{\epsilon}{2}.
    \end{align*}
    It follows that
    $$ R(P) = U_P - L_P < \brac{A + \frac{\epsilon}{2}} - \brac{A - \frac{\epsilon}{2}} = \epsilon. $$
\end{itemize}
\end{solution}
\end{problem}

\begin{problem}
Prove: if $f$ is a monotonic function on the interval $[a,b]$, then $f$ is integrable on $[a,b]$. (You may use any definition of the integral you wish.)
\begin{solution}
By setting $f := -f$ if need be, we may assume without loss of generality that $f$ is increasing. Notice that $f$ is bounded, because $f(a) \leq f(x) \leq f(b)$ for any $x \in [a,b]$. If $f(a) = f(b)$ then $f$ is constant and there is nothing to prove; therefore, we can assume that $f(a) \neq f(b)$. Given a partition $P = \set{a = x_0 < x_1 < \cdots < x_n = b}$, let
\begin{align*}
    L_P &= \sum_{i=1}^n m_i (x_i - x_{i-1}) = \sum_{i=1}^n f(x_{i-1}) (x_i - x_{i-1}), \\
    U_P &= \sum_{i=1}^n M_i (x_i - x_{i-1}) = \sum_{i=1}^n f(x_i) (x_i - x_{i-1}),
\end{align*}
where $m_i = \inf\{f(x) \mid x_{i-1} \leq x \leq x_i\}$, $M_i = \sup\{f(x) \mid x_{i-1} \leq x \leq x_i\}$, and where we used the fact that $f$ is increasing for the rightmost equalities.
Recall that a function $f$ is said to be \emph{Darboux-integrable} if for any $\epsilon > 0$ there exists a $\delta > 0$ such that
$$ U_P - L_P < \epsilon $$
for any partition $P$ such that $\mu(P) < \delta$, where $\mu(P) = \max\{x_i - x_{i-1} \mid i = 1, \ldots, n\}$ is the \emph{mesh} of $P$. Let $\epsilon > 0$, and $\delta = (f(b) - f(a))^{-1} \epsilon$.
\begin{align*}
    U_P - L_p &= \sum_{i=1}^n f(x_i) (x_i - x_{i-1}) - \sum_{i=1}^n f(x_{i-1}) (x_i - x_{i-1}) \\
    &= \sum_{i=1}^n (f(x_i) - f(x_{i-1})) (x_i - x_{i-1}) \\
    &\leq \sum_{i=1}^n (f(x_i) - f(x_{i-1})) \mu(P) \\
    &= (f(b) - f(a)) \mu(P) \\
    &< \epsilon.
\end{align*}
\end{solution}
\end{problem}

\begin{problem}
Define 
    \[f(x) = \int_0^\pi \frac{\sin(xt)}{t} \; dt , \; x\in \R\]
    \begin{itemize}
        \item [(a)] Prove $f$ is a continuous function on $\R$
        \item [(b)] Prove $f$ is differentiable on $\R$ and compute $f'(x)$. 
    \end{itemize}
\begin{solution}
a) Let $c\in \R$ and given $\epsilon > 0$, for any $x$ by the mean value theorem there exists $d_x\in B(c,|x-c|)$ such that $\sin(tx) - \sin(tc)= \left(\frac{d}{dx} \sin(tx)|_{x=d_x}\right) (x-c) = t\cos(td_x)(x-c) $ hence,
    \[\begin{array}{rcl}\displaystyle \left|f(x)-f(c) \right|&=& \displaystyle\left|\int_0^\pi \frac{\sin(xt)}{t} \; dt - \int_0^\pi \frac{\sin(ct)}{t} \; dt \right| \\ \\ 
    &=& \displaystyle\left|\int_0^\pi \frac{\sin(xt)-\sin(ct)}{t}  \; dt \right| 
    \\ \\ 
    &\leq & \displaystyle\int_0^\pi \left| \frac{\sin(xt)-\sin(ct)}{t}  \; \right|  dt 
    \\ \\ 
    &\leq & \displaystyle\int_0^\pi \left| \frac{t \cos(td_x)(x-c)}{t}  \; \right|  dt \leq \pi |x-c| 
    \end{array}\]
    Thus for $\delta < \frac{\varepsilon}{\pi} $ we get $|f(x)-f(c)|< \epsilon$ whenever $|x-c| < \delta$\newline 

    b) For any $x\in \R$ function $g(t) = \begin{cases}
        \frac{\sin(xt)}{t} & t\ne 0 \\ 
        x & t=0\end{cases} $
    is  continuous for $t\in [0,\pi]$ continuous on compact set is uniform continuous so limit of the integral is the integral of the limit. Then 
    \[\begin{array}{rcl}
    \displaystyle \lim_{h\to 0} \frac{f(x+h)-f(x)}{h}&=&\displaystyle \lim_{h\to 0} \frac{1}h \left[\int_0^\pi \frac{\sin(xt+ht)}{t}\; dt - \int_0^\pi \frac{\sin(xt)}{t}\; dt \right]\\ \\ 
    &=&\displaystyle \lim_{h\to 0} \frac{1}h \left[\int_0^\pi \frac{\sin(xt+ht)-\sin(xt)}{t}\; dt  \right]\\ \\ 
    &=&\displaystyle \lim_{h\to 0} \frac{1}h \left[\int_0^\pi \frac{2\cos\left(\frac{xt+ht+xt}{2}\right)\sin\left(\frac{ht}{2}\right)}{t}\; dt  \right]
    \\ \\ 
    &=&\displaystyle  \left[\int_0^\pi \lim_{h\to 0} \frac{1}h \frac{2\cos\left(\frac{xt+ht+xt}{2}\right)\sin\left(\frac{ht}{2}\right)}{t}\; dt  \right] \\ 
    &=& \displaystyle \int_0^\pi \cos(xt)\; dt = \left.\frac{\sin(xt)}{x}\right|_0^\pi = \begin{cases}
        \frac{\sin(\pi x)}{x} & x \ne 0 \\ \pi & x = 0 
    \end{cases}
    \end{array}\]
\end{solution}
\end{problem}

\begin{problem}
For every integer $n\geq 0 $, let $I_n = \int_0^{\frac{\pi}{2}} \cos^n (t) \; dt$. 
    \begin{itemize}
        \item [(a)] Show that $(n+2)I_{n+2} = (n+1) I_n$. Then use it to find explicitly $I_{2n}$ and $I_{2n+1}$
        \item [(b)] Show that $\lim_{n\to \infty} \frac{I_{n+1}}{I_n} = 1 $
        \item [(c)] Show that $\{(n+1)I_nI_{n+1}\}$ is a constant sequence. Then conclude that $\lim_{n\to \infty} I_n \sqrt{2n}= \sqrt{\pi}$
    \end{itemize}
\begin{solution}
a) Using integration by parts 
    \[\begin{array}{rcl}
    \displaystyle I_{n+2} &=&  \displaystyle \int_0^{\pi/2} \underbrace{\cos(t)}_{dv}\underbrace{\cos^{n+1}(t)}_{u}\; dt \\ \\ &=&  \displaystyle \left.\sin(t)\cos^{n+1}(t)\right|_0^{\pi/2} - \int_0^{\pi/2} \sin(t)((n+1)\cos^n(t) (-\sin t))\; dt  \\ \\ 
     &=&  \displaystyle (n+1) \int_0^{\pi/2} \underbrace{\sin^2(t)}_{1-\cos^2(t)}\cos^n(t) \; dt \\ &=& (n+1)[ I_n - I_{n+2}]
    \end{array}\]
    So we get $(n+2)I_{n+2}=(n+1)I_n$\newline 

    By the previous we have $\displaystyle I_{n} = \frac{n-1}{n} I_{n-2}$
    and we have $I_0 = \int_0^{\pi/2} 1 \; dt = \frac{\pi}{2}$ and $I_1 = \int_0^{\pi/2} \cos t \; dt = \sin t|_0^{\pi/2} = 1 $ thus for $n=2k+1$ odd then 
    \[I_{2k+1} = \frac{2k}{2k+1} I_{2k-1} = \frac{2k}{2k+1}\cdot \frac{2k-2}{2k-1}\cdots \frac23  =\frac{2k}{2k+1}\cdot \frac{2k}{2k}\cdot \frac{2k-2}{2k-1}\cdot\frac{2k-2}{2k-2}\cdots \frac23 = \frac{4^k (k!)^2}{(2k+1)!} \]
    for $n=2k$ even 
    \[I_{2k} = \frac{2k-1}{2k} I_{2k-2} = \frac{2k-1}{2k} \cdot \frac{2k-3}{2k-2}\cdots \frac12 \cdot I_0 =\frac{2k}{2k} \cdot\frac{2k-1}{2k} \cdot \frac{2k-2}{2k-2}\cdot \frac{2k-3}{2k-2}\cdots \frac12 \cdot \frac{\pi}2=\frac{(2k)!\pi}{4^k (k!)^2 (2)}  \]

    b) Note that $$I_{n+1} - I_n = \int_0^{\pi/2} \cos^{n+1} (t) - \cos^n(t) \; dt = \int_0^{\pi/2} \cos^n(t) ( \cos t - 1)\; dt\leq 0 $$
    Since $1\geq \cos t, \forall t$ hence $I_{n+1}\leq I_n$ for all $n$ so the sequence $I_n$ is decreasing. Thus 
    $$\begin{array}{rcccl}
    \displaystyle I_{n-1} &\geq& \displaystyle I_n &\geq& I_{n+1} \\ 
     \displaystyle 1 &\geq& \displaystyle \frac{I_n}{I_{n-1}}&\geq& \displaystyle\frac{I_{n+1}}{I_{n-1}} \\ 
     1  &\geq& \frac{I_n}{I_{n-1}} &\geq& \frac{n}{n+1}
        
    \end{array}$$
    Taking the limit by squeeze theorem we get $\lim \frac{I_n}{I_{n-1}}\to 1$\newline 

    c) Multiplying $(n+2)I_{n+2}=(n+1)I_n$ by $I_{n+1}$ gives $(n+2)I_{n+1}I_{n+2}=(n+1)I_nI_{n+1}$, so the sequence is constant. Its value can be read from the odd and even $I_n$: 
    \[(2k+1)I_{2k+1}I_{2k}= (2k+1)\frac{4^k (k!)^2}{(2k+1)!}\frac{(2k)!\pi}{4^k (k!)^2 (2)}= \frac{\pi}{2}\]
    Thus 
    \[ \sqrt{2(n)}I_n \sqrt{2(n+1)}I_{n+1} = \frac{\pi}{2(n+1)}2\sqrt{n(n+1)}= \pi \sqrt{\frac{n}{n+1}}\]
    By b), $\dfrac{\sqrt{2n}I_n}{\sqrt{2(n+1)}I_{n+1}} = \sqrt{\dfrac{n}{n+1}}\,\dfrac{I_n}{I_{n+1}} \to 1$, so
    \[ \left(\sqrt{2n}I_n\right)^2 = \sqrt{2n}I_n \sqrt{2(n+1)}I_{n+1}\cdot \frac{\sqrt{2n}I_n}{\sqrt{2(n+1)}I_{n+1}} \to \pi\cdot 1 = \pi, \]
    and since $I_n>0$, $\lim \sqrt{2n}I_n =\sqrt{\pi}$.
\end{solution}
\end{problem}

\begin{problem}
Suppose $f: [0,1] \to [0, \infty)$ is a continuous function. Let $A = \int_0^1 f(x) \,dx$. Prove that
$$ \sqrt{1 + A^2} \leq \int_0^1 \sqrt{1 + [f(x)]^2} \,dx \leq 1 + A. $$
\begin{solution}
For the second inequality, notice that, for any $x \geq 0$,
$$ 1 + x^2 \leq 1 + 2x + x^2 = (1 + x)^2 \implies \sqrt{1 + x^2} \leq 1 + x. $$
Therefore,
$$ \int_0^1 \sqrt{1 + [f(x)]^2} \,dx \leq \int_0^1 1 + f(x) \,dx = 1 + A. $$
On the other hand, the function $\phi(x) = \sqrt{1 + x^2}$ is convex, because
$$ \phi''(x) = \brac{\frac{x}{\sqrt{1 + x^2}}}' = \frac{\sqrt{1 + x^2} - \frac{x^2}{\sqrt{1 + x^2}}}{1 + x^2} = \frac{1}{(1 + x^2)^{3/2}} > 0 $$
for all $x \in \R$. By Jensen's inequality,
$$ \sqrt{1 + A^2} = \phi\brac{\int_0^1 f(x) \,dx} \leq \int_0^1 \phi(f(x)) \,dx = \int_0^1 \sqrt{1 + [f(x)]^2} \,dx. $$

Alternative proof for the first inequality: Notice that the graph of $F(x) = \int_0^x f(t) \,dt$ passes through both $(0,0)$ and $(1,A)$. The left-hand side is the straight line distance between the two points, while the middle part is the arc length of the graph of $F$ (which then has to be at least as long.)
\end{solution}
\end{problem}

\begin{problem}
Suppose $f$ is a real, continuously differentiable function on $[0,1]$, $f(0)=f(1)=0$ and \[\int_0^1 f(x) e^x \; dx = 1\]
    Prove 
    \begin{itemize}
        \item [(a)] $\int_0^1 f'(x) e^x \; dx = -1 $ and 
        \item[(b)] $\int_0^1 (f'(x))^2 \; dx \geq \frac{2}{e^2}$.
    \end{itemize}
\begin{solution}
Integrating by parts we get 
    \[\int_0^1 f(x) e^x\; dx = e^x f(x)|_0^1 - \int_0^1 f'(x) e^x \; dx \Rightarrow 1 = - \int_0^1 f'(x) e^x\; dx \Rightarrow \int_0^1 f'(x)e^x\; dx = -1 \]
    For the second part, apply the Cauchy Schwartz integral inequality
    \[\begin{array}{rcl}
    \dd \int_0^1 (f'(x))^2 \; dx\int_0^1 (e^x)^2 \; dx &\geq& \dd \left( \int_0^1 f'(x) e^x\; dx\right)^2 \\ 
    \dd \int_0^1 (f'(x))^2 \cdot \left( \frac{e^2}{2} - \frac{1}{2}\right) &\geq&1 \\ 
    \dd \int_0^1 (f'(x))^2 \; dx  &\geq& \dd \frac{2}{e^2-1} 
    \\ 
    \dd \int_0^1 (f'(x))^2 \; dx  &\geq& \dd \frac{2}{e^2} 
    \end{array}\]
\end{solution}
\end{problem}
\begin{rmkbox}{Remark}
In the version of this problem we worked from, part (a) asked to show $\int_0^1 f'(x) e^x \, dx = 1$. Integration by parts gives $-1$, as the solution shows; part (b) is not affected, since the integral is squared.
\end{rmkbox}

\begin{problem}
Suppose that $f$ is a real, continuously differentiable function on $[a,b]$ with $b > a \ge 0$, $f(a) = f(b) = 0$ and 
\begin{align*}
    \int_a^b f(x)^2 \,dx = 1.
\end{align*}
Prove that 
\begin{align*}
    \int_a^b x f(x) f'(x) \,dx = -\dfrac{1}{2}
\end{align*}
and that 
\begin{align*}
    \int_a^b x (f'(x))^2 \,dx \cdot \int_a^b x (f(x))^2 \,dx > \dfrac{1}{4}.
\end{align*}
Be aware that the second inequality is a strict inequality.
\begin{solution}
Use integral by parts and note that $f(a)=f(b)=0$, we have
\begin{align*}
    \int_a^bxf(x)f'(x)dx&=xf(x)f(x)|_a^b-\int_a^b(xf(x))'f(x)dx\\
    &=-\int_a^b(f(x)+xf'(x))f(x)dx\\
    &=-\int_a^bxf(x)f'(x)dx -\int_a^bf(x)^2dx 
\end{align*} and thus 
\begin{align*}
    \int_a^bxf(x)f'(x)dx=-\dfrac{1}{2} \int_a^bf(x)^2dx=-\dfrac{1}{2}.
\end{align*}
Note that 
\begin{align*}
    \brac{\int_a^b g(x)h(x)dx}^2\le \int_a^b g(x)^2dx\int_a^bh(x)^2dx,
\end{align*}
we have 
\begin{align*}
    \dfrac{1}{4}&=\brac{\int_a^b xf(x)f'(x)dx}^2\le \int_a^b x(f(x))^2dx \int_a^b x(f'(x))^2dx.
\end{align*}
The equality holds only if there is some $\lambda\ne 0$ such that 
\begin{align*}
    \sqrt{x} f(x)=\lambda\sqrt{x}f'(x),\text{ for all }x\in [a,b],
\end{align*}
and thus 
\begin{align*}
    f(x)=\lambda f'(x),\text{ for all }x\in (a,b).
\end{align*}
Then there is some $C\in\mathbb{R}$ such that 
\begin{align*}
    f(x)=Ce^{x/\lambda},\text{ for all }x\in (a,b).
\end{align*}
Clearly, $C \neq 0$ since the integral is non-zero. Then this function is never 0, which contradicts the assumption that $f(a) = f(b) = 0$. Therefore we must have a strict inequality.
\end{solution}
\end{problem}
\begin{rmkbox}{Remark}
In the version of this problem we worked from, the hypothesis was $\int_a^b f(x)\,dx = 1$. The first identity needs $\int_a^b f(x)^2\,dx = 1$, as stated above.
\end{rmkbox}
